% GENERATED FILE. Edit canonical lessons, then run npm run generate:books.
% canonical-source-hash: fnv1a64:e0c4f322c01755fa
% canonical-lessons: UR-C546-mashaqqat, UR-C546-hath-batana, UR-C546-asani, UR-C546-jujhna, UR-C546-paimana

\chapter{The Hard Way}
\label{ch:ur-the-hard-way}
\hlchaptermodality{546}

\begin{chapteropening}
\textbf{By the end of this chapter:} \emph{I can talk about hard work and what makes things easier.}

Complete the last lesson of chapter 546: I can talk about hard work and what makes things easier.
\end{chapteropening}

\section[\texorpdfstring{\ur{مشقت} (mashaqqat)}{مشقت (mashaqqat)}]{\ur{مشقت} (mashaqqat) — toil; hard labour}

\label{lesson:UR-C546-mashaqqat}



\begin{quote}
Before the new one: say the Urdu for a stage, then the Urdu for durable.
\end{quote}

\subsection*{You'll want to know: \ur{مشقت}}
\textbf{\ur{مشقت}} — \emph{mashaqqat} — \textquotedblleft{}toil; hard labour\textquotedblright{}.

\begin{quote}
\ur{اگر} \ur{مشقت} \ur{نہ} \ur{ہوتی} \ur{تو} \ur{یہ} \ur{پل} \ur{کبھی} \ur{نہ} \ur{بنتا۔}
\end{quote}

\begin{quote}
\emph{agar mashaqqat na hotī to yih pul kabhī na bantā.}
\end{quote}

\begin{quote}
Without hard labour this bridge would not have been built at all.
\end{quote}

Feminine.

\begin{cousinweb}
\textbf{\ur{مشقت}} (\emph{mashaqqat}) is usually traced, through Persian, to Arabic \emph{mashaqqa}, \textquotedblleft{}hardship, toil\textquotedblright{}.
\end{cousinweb}

\begin{grammarlens}[title={What you've built}]
One more word for saying what happens if.
\end{grammarlens}

\subsection*{Your turn}
\begin{itemize}
\raggedright
  \item \emph{Say it:} \emph{mashaqqat}
  \item \emph{Say it:} \emph{mashaqqat}, once more
  \item \emph{Say it:} \emph{agar mashaqqat na hotī to yih pul kabhī na bantā.} — Without hard labour this bridge would not have been built at all.
  \item \emph{Recall:} say the Urdu for a stage, then the Urdu for durable, then say \emph{mashaqqat} again
\end{itemize}

\begin{tcolorbox}[breakable,colback=teal!4,colframe=teal!35!black,title={Before you move on}]
What does \ur{مشقت} mean? (\textquotedblleft{}Toil; hard labour\textquotedblright{}.) Say it once more.
\end{tcolorbox}
\section[\texorpdfstring{\ur{ہاتھ} \ur{بٹانا} (hāth baṭānā)}{ہاتھ بٹانا (hāth baṭānā)}]{\ur{ہاتھ} \ur{بٹانا} (hāth baṭānā) — to lend a hand}

\label{lesson:UR-C546-hath-batana}



\begin{quote}
Before the new one: say the Urdu for durable, then the Urdu for toil.
\end{quote}

\subsection*{You'll want to know: \ur{ہاتھ} \ur{بٹانا}}
\textbf{\ur{ہاتھ} \ur{بٹانا}} — \emph{hāth baṭānā} — \textquotedblleft{}to lend a hand\textquotedblright{}.

\begin{quote}
\ur{اگر} \ur{سب} \ur{نے} \ur{ہاتھ} \ur{بٹایا} \ur{تو} \ur{کام} \ur{ختم} \ur{ہو} \ur{گا۔}
\end{quote}

\begin{quote}
\emph{agar sab ne hāth baṭāyā to kām khatm ho gā.}
\end{quote}

\begin{quote}
If everyone lends a hand, the work will get finished.
\end{quote}

Past, with \ur{نے} (he or she …): \ur{اس} \ur{نے} \ur{امی} \ur{کا} \ur{ہاتھ} \ur{بٹایا}, \emph{us ne ammī kā hāth baṭāyā}.

\begin{cousinweb}
\textbf{\ur{ہاتھ} \ur{بٹانا}} (\emph{hāth baṭānā}) is built from \textbf{\ur{ہاتھ}} (\emph{hāth}), \textquotedblleft{}a hand\textquotedblright{}, and \emph{baṭānā}, \textquotedblleft{}to share in\textquotedblright{}, which sits beside the Hindi cousin \emph{bā̃ṭnā}, \textquotedblleft{}shares\textquotedblright{}, that CDIAL lists under Sanskrit \emph{váṇṭati} (CDIAL 11238): to take a share of someone's work.
\end{cousinweb}

\begin{grammarlens}[title={What you've built}]
One more word for saying what happens if.
\end{grammarlens}

\subsection*{Your turn}
\begin{itemize}
\raggedright
  \item \emph{Say it:} \emph{hāth baṭānā}
  \item \emph{Say it:} \emph{hāth baṭānā}, once more
  \item \emph{Say it:} \emph{agar sab ne hāth baṭāyā to kām khatm ho gā.} — If everyone lends a hand, the work will get finished.
  \item \emph{Recall:} say the Urdu for durable, then the Urdu for toil, then say \emph{hāth baṭānā} again
\end{itemize}

\begin{tcolorbox}[breakable,colback=teal!4,colframe=teal!35!black,title={Before you move on}]
What does \ur{ہاتھ} \ur{بٹانا} mean? (\textquotedblleft{}To lend a hand\textquotedblright{}.) Say it once more.
\end{tcolorbox}
\section[\texorpdfstring{\ur{آسانی} (āsānī)}{آسانی (āsānī)}]{\ur{آسانی} (āsānī) — ease}

\label{lesson:UR-C546-asani}



\begin{quote}
Before the new one: say the Urdu for toil, then the Urdu for to lend a hand.
\end{quote}

\subsection*{You'll want to know: \ur{آسانی}}
\textbf{\ur{آسانی}} — \emph{āsānī} — \textquotedblleft{}ease\textquotedblright{}.

\begin{quote}
\ur{اگر} \ur{نقشہ} \ur{ہوتا} \ur{تو} \ur{ہم} \ur{آسانی} \ur{سے} \ur{پہنچ} \ur{جاتے۔}
\end{quote}

\begin{quote}
\emph{agar naqshā hotā to ham āsānī se pahuṅch jāte.}
\end{quote}

\begin{quote}
If we had a map, we would get there easily.
\end{quote}

Feminine.

\begin{cousinweb}
\textbf{\ur{آسانی}} (\emph{āsānī}) is usually traced to Persian \emph{āsānī}, \textquotedblleft{}ease\textquotedblright{}, built on \emph{āsān}, \textquotedblleft{}easy\textquotedblright{}.
\end{cousinweb}

\begin{grammarlens}[title={What you've built}]
One more word for saying what happens if.
\end{grammarlens}

\subsection*{Your turn}
\begin{itemize}
\raggedright
  \item \emph{Say it:} \emph{āsānī}
  \item \emph{Say it:} \emph{āsānī}, once more
  \item \emph{Say it:} \emph{agar naqshā hotā to ham āsānī se pahuṅch jāte.} — If we had a map, we would get there easily.
  \item \emph{Recall:} say the Urdu for toil, then the Urdu for to lend a hand, then say \emph{āsānī} again
\end{itemize}

\begin{tcolorbox}[breakable,colback=teal!4,colframe=teal!35!black,title={Before you move on}]
What does \ur{آسانی} mean? (\textquotedblleft{}Ease\textquotedblright{}.) Say it once more.
\end{tcolorbox}
\section[\texorpdfstring{\ur{جوجھنا} (jūjhnā)}{جوجھنا (jūjhnā)}]{\ur{جوجھنا} (jūjhnā) — to grapple with (a hard task)}

\label{lesson:UR-C546-jujhna}



\begin{quote}
Before the new one: say the Urdu for to lend a hand, then the Urdu for ease.
\end{quote}

\subsection*{You'll want to know: \ur{جوجھنا}}
\textbf{\ur{جوجھنا}} — \emph{jūjhnā} — \textquotedblleft{}to grapple with (a hard task)\textquotedblright{}.

\begin{quote}
\ur{اگر} \ur{ہم} \ur{ساتھ} \ur{جوجھے} \ur{تو} \ur{یہ} \ur{مشکل} \ur{ختم} \ur{ہو} \ur{گی۔}
\end{quote}

\begin{quote}
\emph{agar ham sāth jūjhe to yih mushkil khatm ho gī.}
\end{quote}

\begin{quote}
If we grapple with it together, this trouble will end.
\end{quote}

Past (he …): \ur{وہ} \ur{بیماری} \ur{سے} \ur{جوجھا}, \emph{voh bīmārī se jūjhā}.

\begin{cousinweb}
\textbf{\ur{جوجھنا}} (\emph{jūjhnā}) is inherited: CDIAL lists the Hindi cousin \emph{jūjhnā}, \textquotedblleft{}to fight\textquotedblright{}, under Sanskrit \emph{yudhyatē}, \textquotedblleft{}fights\textquotedblright{} (CDIAL 10502).
\end{cousinweb}

\begin{grammarlens}[title={What you've built}]
One more word for saying what happens if.
\end{grammarlens}

\subsection*{Your turn}
\begin{itemize}
\raggedright
  \item \emph{Say it:} \emph{jūjhnā}
  \item \emph{Say it:} \emph{jūjhnā}, once more
  \item \emph{Say it:} \emph{agar ham sāth jūjhe to yih mushkil khatm ho gī.} — If we grapple with it together, this trouble will end.
  \item \emph{Recall:} say the Urdu for to lend a hand, then the Urdu for ease, then say \emph{jūjhnā} again
\end{itemize}

\begin{tcolorbox}[breakable,colback=teal!4,colframe=teal!35!black,title={Before you move on}]
What does \ur{جوجھنا} mean? (\textquotedblleft{}To grapple with (a hard task)\textquotedblright{}.) Say it once more.
\end{tcolorbox}
\section[\texorpdfstring{\ur{پیمانہ} (paimānā)}{پیمانہ (paimānā)}]{\ur{پیمانہ} (paimānā) — a yardstick; a measure}

\label{lesson:UR-C546-paimana}



\begin{quote}
Before the new one: say the Urdu for ease, then the Urdu for to grapple with (a hard task).
\end{quote}

\subsection*{You'll want to know: \ur{پیمانہ}}
\textbf{\ur{پیمانہ}} — \emph{paimānā} — \textquotedblleft{}a yardstick; a measure\textquotedblright{}.

\begin{quote}
\ur{اگر} \ur{پیسہ} \ur{ہی} \ur{پیمانہ} \ur{ہوتا} \ur{تو} \ur{استاد} \ur{امیر} \ur{ہوتے۔}
\end{quote}

\begin{quote}
\emph{agar paisā hī paimānā hotā to ustād amīr hote.}
\end{quote}

\begin{quote}
If money were the only yardstick, teachers would be rich.
\end{quote}

Masculine.

\begin{cousinweb}
\textbf{\ur{پیمانہ}} (\emph{paimānā}) is usually traced to Persian \emph{paymāna}, \textquotedblleft{}a measure; a measuring cup\textquotedblright{}.
\end{cousinweb}

\begin{grammarlens}[title={What you've built}]
One more word for saying what happens if.
\end{grammarlens}

\subsection*{Your turn}
\begin{itemize}
\raggedright
  \item \emph{Say it:} \emph{paimānā}
  \item \emph{Say it:} \emph{paimānā}, once more
  \item \emph{Say it:} \emph{agar paisā hī paimānā hotā to ustād amīr hote.} — If money were the only yardstick, teachers would be rich.
  \item \emph{Recall:} say the Urdu for ease, then the Urdu for to grapple with (a hard task), then say \emph{paimānā} again
\end{itemize}

\begin{tcolorbox}[breakable,colback=teal!4,colframe=teal!35!black,title={Before you move on}]
What does \ur{پیمانہ} mean? (\textquotedblleft{}A yardstick; a measure\textquotedblright{}.) Say it once more.
\end{tcolorbox}
